% !TeX root = ../../quotients-and-descent.tex
\section{The restricted principal census and Jacobian question}\label{q:jacobians}
This section supplies the order-and-type result invoked for the source
factor. Its arithmetic proof is independent of the quotient and descent
sections. The theta step has an additional execution and enclosure
contract; the exact conditional statement is retained below.

\begin{theorem}[Complete simple principal classification with the selected order]
\label{pp:thm:tenprincipal}
There are exactly ten origin-preserving polarized isomorphism classes of
simple principally polarized complex abelian sixfolds $X$ such that
$\operatorname{End}^0(X)\simeq M$ and, under this identification,
$R\subseteq\operatorname{End}(X)$. The containment condition is independent
of the field identification. Eight classes have endomorphism ring $R$
and polarized automorphism group $\mu_6$; two have endomorphism ring
$\OM$ and polarized automorphism group $\mu_{42}$. They belong to four
unpolarized isogeny classes, with the following counts at one CM-type
representative from each orbit:
\[
\begin{array}{c|c|ccc}
\Phi&\text{flip parity from }\Phi_0 &[R]&[\mathfrak a]&[\OM]\\ \hline
\{1,5,11,13,17,19\}&0&2&2&0\\
\{1,5,11,13,19,25\}&1&0&0&1\\
\{1,5,11,17,23,29\}&0&2&2&0\\
\{1,5,11,23,25,29\}&1&0&0&1
\end{array}
\]
In particular, the four specified polarized B/C objects exhaust the
objects in their unpolarized isogeny class satisfying this $R$-containment
condition. No assertion is made here for other endomorphism orders or
for the induced A type.
\end{theorem}
\begin{proof}
Every automorphism of $M$ preserves $\OL$ and takes $1-z^{14}$ to
itself or an associate. Hence it preserves $R=\OL+(1-z^{14})\OM$.
After choosing a field identification, $H_1(X,\Ql)$ has dimension one
over $M$, and its integral lattice is a full fractional $R$-ideal.
Conversely, the positive unimodular forms of
Proposition~\externalref{packet}{pp:prop:markedpolarizations} give abelian varieties
by the Riemann criterion.

We check full endomorphisms and simplicity, not merely a centralizer.
For a CM type $\Phi$, put
$H=\{\sigma\in\operatorname{Gal}(M/\Ql):\sigma\Phi=\Phi\}$.
The splitting of
$\operatorname{End}_{\Ql}(M)=\bigoplus_\sigma M\sigma$
over $\mathbb C$ has disjoint permutation-matrix supports for distinct
$\sigma$. A nonzero coefficient in $M$ has nonzero value at every
embedding, so a term crossing the two Hodge blocks cannot cancel.
Thus the full rational Hodge endomorphism algebra is
$M\mathbin{\rtimes}H=\operatorname{End}_{M^H}(M)$.
The rational weight-one equivalence
\cite[Corollary~6.9]{MilneShimura2017} and Poincar\'e complete
reducibility show that the abelian variety is simple precisely when
$H$ is trivial: otherwise this matrix algebra has nontrivial
idempotents, whereas for trivial $H$ it is the field $M$.

The exact census of the $64$ types gives precisely the four
trivial-stabilizer orbits in the table, each of size $12$.
The two retained implementations are
\path{experiments/fixed_field_cm_census_v98/primary.py} and
\path{experiments/fixed_field_cm_census_v98/independent.py}.
The bounded table derivation
\path{experiments/simple_cm_unmarked_v145/derive.py}
binds their accepted outputs, checks orbit coverage and parity, and
applies Proposition~\externalref{packet}{pp:prop:markedpolarizations}; it does not
re-prove the ideal classification or the unit-norm theorem.

Every polarization on a simple object here has Rosati involution
equal to conjugation on $M$. Indeed its positive trace form extends
to $M\otimes\mathbb R\simeq\mathbb C^6$. An involution exchanging
factors would give trace norm zero on a factor idempotent. On a fixed
complex factor, the identity would give negative trace norm on $i$.
Positivity therefore forces conjugation on every factor. Thus the
marked classification covers every principal polarization.

Fix one primitive type representative. An unmarked polarized
isomorphism conjugates the full endomorphism fields and hence is
semilinear for some $\sigma\in\operatorname{Gal}(M/\Ql)$.
Holomorphicity at this fixed type forces $\sigma\Phi=\Phi$,
so $\sigma=1$. Forgetting the marking therefore identifies no
additional classes or automorphisms at the chosen representative.
Every object in its orbit can be remarked to this representative,
because $R$ is Galois-stable. The table and automorphism assertions
now follow from Proposition~\externalref{packet}{pp:prop:markedpolarizations}.

An unpolarized isogeny likewise induces a field automorphism relating
the types, so distinct primitive orbits are not isogenous. Within one
orbit, after remarking, full fractional lattices in $M$ are
commensurable; an integer multiple of the identity gives an isogeny.
Finally $\sigma_{31}(\Phi_C)=\Phi_B$ places the four distinct
objects of Proposition~\externalref{packet}{pp:prop:actualpolarizations} in the first
orbit. Its count of four proves the restricted exhaustion assertion.
\end{proof}

\begin{corollary}[Jacobian membership in the ten-class list]
\label{pp:cor:jacobianlist}
Of the ten classes in Theorem~\ref{pp:thm:tenprincipal}, the four in
the third row and the single class in the second row are not
canonically polarized Jacobians of smooth curves. The single class
in the fourth row is the known Jacobian of the smooth projective
model of $y^7=x^3-1$.
These signature and automorphism criteria alone leave the four B/C
classes in the first row undecided; Theorem~\ref{pp:thm:jacobiancomplete}
will exclude them using the invariant theta divisor.
If any is a Jacobian, its curve is nonhyperelliptic, has full
automorphism group $C_3$, and is a cyclic triple cover of
$\mathbb P^1$ with eight branch points, four of each nonzero
inertia exponent.
\end{corollary}
\begin{proof}
Strong Torelli \cite[Theorem~12.1]{MilneJacobians2021} realizes an
odd-order polarized Jacobian automorphism $\beta$ by a curve
automorphism up to sign. Squaring the latter induces $\beta^2$
and has the same odd order, by faithfulness on the Jacobian.
For $\beta=\zeta_3$, there are no invariant differentials, so the
curve quotient has genus zero. Riemann--Hurwitz gives eight fixed
points. If $n_1,n_2$ count the two inertia exponents, then
$n_1+n_2=8$ and $n_1+2n_2\equiv0\pmod3$.
The cyclic-cover eigenspace formula
\cite[equation~(2.2)]{LiMantovanPries2021} gives only signatures
$(4,2),(3,3),(2,4)$, respectively for $n_2=1,4,7$.
The third row has signature $(1,5)$ and is therefore excluded.
This is also directly Zarhin's established criterion
$g+2<3|a-b|$ \cite[Theorem~1.1]{Zarhin2007}, here $8<12$;
the obstruction method is not new.
The first row has signature $(3,3)$ and hence requires $n_1=n_2=4$.
A hyperelliptic curve would induce a nontrivial order-three map on
its hyperelliptic base $\mathbb P^1$, with only two fixed points
and at most four curve points above them. This contradicts eight
fixed points. Strong Torelli and the full polarized automorphism
group $\mu_6$ then give full curve automorphism group $C_3$.

For either maximal-order object, use instead $\beta=\zeta_{21}$.
Its eigencharacters are primitive, so again the curve quotient
has genus zero. If $a,b,c$ count branch points with ramification
indices $3,7,21$, Riemann--Hurwitz gives
$7a+9b+10c=26$. Its only nonnegative solution is $(1,1,1)$.
After fixing these three branch points at $0,1,\infty$, the
monodromy residues are $(7u,3v,-7u-3v)$, with
$u\in(\mathbb Z/3)^\times$ and $v\in(\mathbb Z/7)^\times$.
The Chinese remainder theorem makes these twelve tuples a single
orbit under a change of cyclic generator. Kummer theory thus gives
one curve up to complex isomorphism, namely $y^7=x^3-1$.
Its known simple maximal-order CM type is
$\{1,2,4,8,10,16\}$ modulo $21$
\cite[Proposition~6.11, Lemma~6.12 and Proposition~6.14]{LiMantovanPries2021}.
Its odd lifts modulo $42$ are $\{1,23,25,29,31,37\}$, which are
$23$ times the fourth row's representative. That row is therefore
the known Jacobian. The distinct second-row object cannot also be
a Jacobian. The finite ramification and type arithmetic is retained
in \path{experiments/simple_cm_unmarked_v145/jacobian_check.py};
the proof uses the geometric inputs just named and makes no
assertion about isogeny to a Jacobian.
\end{proof}

\subsection{An invariant theta obstruction and the Jacobian classification}
\label{pp:sec:thetacertificate}
Write $B_0,B_1,C_0,C_1$ for the four principal objects of
Proposition~\externalref{packet}{pp:prop:actualpolarizations}, where the subscript records
the prescribed polarization or its $4-t^2$ twist. Throughout this
subsection, Jacobian means the canonical principal polarization of a
smooth curve, up to origin-preserving polarized isomorphism.

\begin{proposition}[The balanced triple-cover count]
\label{pp:prop:trigonaltheta}
Let $C$ be a smooth genus-six cyclic triple cover of $\mathbb P^1$
with eight branch points, four of each nonzero inertia exponent.
If $\delta$ is a deck generator and $G=\ker(1-\delta)\subset J(C)$,
the unique $\delta$-invariant symmetric theta divisor has exactly
$213$ points of $G$ outside its support. Its multiplicities on $G$
are $0,1,2$ with respective frequencies $213,456,60$.
\end{proposition}
\begin{proof}
Move infinity away from the branch locus and write
$y^3=\prod_{i=1}^8(x-t_i)^{a_i}$, with
$a=(1,1,1,1,2,2,2,2)$. Let $P_i$ be the ramification points and
$H$ the fiber at infinity. Then $3P_i\sim H$ and
$D=\sum_iP_i-H$ has degree five and $2D\sim K_C$.
It is the unique invariant theta characteristic: the characteristic
torsor has translation group $J(C)[2]$, on which $1-\delta$ is
invertible since $(1-\delta)(2+\delta)=3$.

Put $W=\{b\in\mathbb F_3^8:\sum_i b_i=0\}$.
For representatives $b_i\in\{0,1,2\}$, define
\[
 E_b=\sum_i b_iP_i-\frac{\sum_i b_i}{3}H.
\]
The map $b\mapsto[E_b]$ takes values in $G$ and has kernel
$\langle a\rangle$. Indeed $a$ gives $\operatorname{div}(y)$;
conversely, a function with invariant principal divisor is a
$\delta$-eigenfunction $y^j h(x)$, and its orders at the $P_i$
force $b_i=ja_i$ modulo three. Both $W/\langle a\rangle$ and
$G$ have order $3^6$, the latter by the eigenvalues of $1-\delta$
on $H^1(C,\mathbb C)$. This identifies every point of $G$.

For $m=1+\sum_i b_i/3$, the three eigensheaves for sections of
$\mathcal O_C(D+E_b)$ on $\mathbb P^1$ have degrees
\[
 e_k=-m-4k+\sum_i\left\lfloor\frac{ka_i+1+b_i}{3}\right\rfloor,
 \qquad k=0,1,2.
\]
This follows by taking orders of $y^k h(x)$ at the branch points
and infinity. Thus $h^0(D+E_b)=\sum_k\max(e_k+1,0)$, independently
of the distinct branch positions. The exact finite count on the
$729$ classes gives sorted degree triples
\[
\begin{array}{c|rrrrr}
(e_0,e_1,e_2)_{\rm sorted}&(-1,-1,-1)&(-2,-1,0)&(-2,-2,1)&(-3,0,0)&(-3,-1,1)\\
\hline
\text{frequency}&213&456&24&24&12.
\end{array}
\]
The retained finite certificate is specified in
\path{experiments/bc_jacobian_v147/adoption.json}; its domain is
exactly the quotient just proved complete. Riemann's singularity
theorem identifies these dimensions with theta multiplicities
\cite[Introduction, smooth-curve statement]{CasalainaMartinKass2012}.
Translation by $D$ identifies the theta support in $\operatorname{Pic}^0(C)$
with $h^0(D+E_b)>0$. The principal line bundle has $h^0=1$,
so its invariant class determines a unique effective divisor.
The displayed frequencies give the assertion.
\end{proof}

\begin{theorem}[Jacobian classification under the enclosure contract]
\label{pp:thm:jacobiancomplete}
Assume the trusted Arb/acb enclosure and execution contract specified
below for the retained theta certificates. Exactly one of the ten
classes in Theorem~\ref{pp:thm:tenprincipal} is a canonically polarized
Jacobian of a smooth curve: the known class of the smooth projective
model of $y^7=x^3-1$. In particular, none of $B_0,B_1,C_0,C_1$ is
such a Jacobian.
\end{theorem}
\begin{proof}
For each model, $\delta=\zeta_3$ preserves the principal polarization
and satisfies $\delta^2+\delta+1=0$. Symmetric representatives form
a torsor under two-torsion, so the same invertibility argument gives
a unique invariant line-bundle class. A principal ample line bundle
has a one-dimensional space of sections, hence determines a unique
effective divisor in this class. Its support is stable under negation.
If a B/C model were a Jacobian, the necessary geometry in
Corollary~\ref{pp:cor:jacobianlist} would identify this divisor with
the one in Proposition~\ref{pp:prop:trigonaltheta}. Inverting the
deck generator does not change its fixed kernel or invariant divisor.

The retained certificates give, on each model separately, $107$
distinct nonzero classes modulo sign in $\ker(1-\delta)$ at which
the theta section is nonzero. The exact-coordinate sets are the
indexed representatives $0,\ldots,106$ in the retained pair map.
Their counts are
\[
\begin{array}{c|rrrr}
\text{model}&B_0&B_1&C_0&C_1\\\hline
\text{nonzero pairs}&107&107&107&107\\
\text{distinct points}&214&214&214&214.
\end{array}
\]
The kernel is killed by three, so a nonzero point differs from its
negative. Each model therefore supplies $214>213$ points outside
the invariant divisor, contradicting the Jacobian count. No origin
value or sum across different models is used. Combining these four
exclusions with Corollary~\ref{pp:cor:jacobianlist} proves the tally.
\end{proof}

The analytic computation starts from the exact ideal lattice, actual
CM embeddings and principal alternating form. It verifies an integral
symplectic basis, recomputes the $\delta$ action and its unique
invariant quadratic refinement in that basis, and constructs interval
periods. For normalized lattice $[I,\tau]$, a fixed point has the
form $z=(n+\tau m)/3$ with $(1-F)(n,m)^t=0$ modulo three.
The half-characteristic is read from that same refinement. These
conventions and their geometric identification are reviewed in
\path{experiments/bc_jacobian_v147/PERIOD_SETUP_REVIEW.md}.

The nonvanishing records are exact dyadic rectangles for squared theta
values. An independent rational check verifies that a real or imaginary
interval excludes zero and that all representatives are disjoint modulo
sign. This check does not reimplement the analytic evaluator: the
conclusion assumes correctness of the pinned Arb/acb routines, the
stated enclosure/serialization interface, and the admitted execution
and runtime provenance. An enclosure excluding zero for the square
implies nonvanishing of the section; it does not determine a zero's
multiplicity. The cumulative certificate is
\path{experiments/bc_jacobian_v147/certificate_run07b_adoption.json},
which binds the eight accepted retained packages; the geometric consequence
is bound by \path{experiments/bc_jacobian_v147/jacobian_classification_adoption.json}.
The separate hostile review in
\path{experiments/bc_jacobian_v147/session_v150/CLASSIFICATION_REVIEW03.md}
checks all $428$ raw rectangles and the complete implication above, including
the two explicitly admitted runtime variants. It does not prove the analytic
library implementation or its equivalence to source.

The theorem concerns the specified endomorphism-order category and
polarized isomorphism. It does not exclude isogenies to Jacobians or
address arbitrary CM orders, and establishes no additional Hodge
algebraicity or Hilbert-modular orbit membership.

